%%
%% Track 1 — Position Paper Template
%% CHIWork 2026 Workshop:
%%   "Interrogating GenAI Augmentation for CHIworkers:
%%    Strategies for Professional Autonomy and Accountability"
%%
%% Workshop website: https://haukecornell.github.io/chiwork-interrogating-aiwork-workshop-template/
%% Workshop day: June 22, 2026 — Linz, Austria
%%
%% INSTRUCTIONS
%% ─────────────────────────────────────────────────────────────
%% - Maximum 4 pages including references
%% - Use this template as-is; do not change fonts, margins, or class options
%% - After acceptance, submit to arXiv and add the arXiv ID to the
%%   \workshopNote command below before your final upload
%% - Submit the PDF via the workshop Google Form (link on website)
%%
%% TOPICS (position papers should engage with at least one):
%%   • AI dependency & de-skilling in HCI work
%%   • Accountability & ownership in AI-augmented workflows
%%   • Ethical commitments, empathy, and co-thinking
%%   • Responsible practice, disclosure, and professional boundaries
%%   • Critique of specific AI tools in HCI/design/research contexts
%% ─────────────────────────────────────────────────────────────

\documentclass[sigconf]{acmart}

%% Strip ACM rights block (not applicable for workshop arXiv papers)
\settopmatter{printacmref=false}
\renewcommand\footnotetextcopyrightpermission[1]{}
\pagestyle{plain}

%% Workshop conference metadata
\acmConference[CHIWork '26]{CHIWork '26 Workshop: Interrogating GenAI Augmentation for CHIworkers}{June 22, 2026}{Linz, Austria}
\acmYear{2026}

%% ── Workshop identification note ──────────────────────────────
%% This command prints a note identifying this as a workshop submission.
%% Before uploading to arXiv, replace "arXiv:XXXX.XXXXX" with your ID.
\newcommand{\workshopNote}{%
  \begin{quote}
    \small\itshape
    This position paper is submitted to the CHIWork 2026 Workshop:
    \emph{Interrogating GenAI Augmentation for CHIworkers:
    Strategies for Professional Autonomy and Accountability}
    (June 22, 2026, Linz, Austria).
    Workshop proposal: \cite{sandhaus2026interrogating}.
    %% Uncomment and fill in after arXiv upload:
    %% arXiv preprint: arXiv:XXXX.XXXXX
  \end{quote}%
}

%% ── Your paper metadata ───────────────────────────────────────
\title{Your Position Paper Title}

%% Repeat \author{} block for each author
\author{First Author}
\affiliation{%
  \institution{Your Institution}
  \city{City}
  \country{Country}}
\email{author@example.com}

\author{Second Author}
\affiliation{%
  \institution{Second Institution}
  \city{City}
  \country{Country}}
\email{author2@example.com}

\renewcommand{\shortauthors}{Author et al.}

%% ── Abstract (max ~150 words) ────────────────────────────────
\begin{abstract}
  Replace this with your abstract. Briefly state your position,
  the problem or tension you address, and your key argument or
  contribution. Aim for 100--150 words.
\end{abstract}

%% ── Keywords ─────────────────────────────────────────────────
\keywords{Generative AI, accountability, HCI practice, professional ethics}
%% Add or replace with your own keywords

%% ─────────────────────────────────────────────────────────────
\begin{document}
\maketitle

%% Print the workshop identification note (required)
\workshopNote

%% ─────────────────────────────────────────────────────────────
\section{Introduction}
%% ─────────────────────────────────────────────────────────────

State the problem or tension your paper addresses. What aspect of
GenAI use in HCI workflows concerns you? What is the accountability
or autonomy gap you have observed in your own practice or research?

Briefly introduce your position and the main argument you will make.
This section should be no more than half a page.


%% ─────────────────────────────────────────────────────────────
\section{Background and Motivation}
%% ─────────────────────────────────────────────────────────────

Situate your position in relevant prior work. You do not need a
comprehensive literature review—this is a position paper, not a full
paper. Cite the work that most directly motivates your argument.

You might engage with themes from the workshop such as:
\begin{itemize}
  \item AI dependency and de-skilling \cite{sandhaus2026interrogating}
  \item The accountability gap in AI-augmented workflows
  \item Tensions between efficiency and deep work
  \item Soft resistance and alternative practices
\end{itemize}


%% ─────────────────────────────────────────────────────────────
\section{Position / Argument}
%% ─────────────────────────────────────────────────────────────

State your position clearly. This is the core of the paper.

You may structure this section as:
\begin{itemize}
  \item A claim about responsible/irresponsible AI use in a specific context
  \item A critique of a current norm or tool
  \item A proposed framework, heuristic, or practice
  \item An empirical observation from your own work or research
\end{itemize}

Be concrete. Use specific examples from practice—what did you observe,
experience, or study? What makes it an instance of the themes
this workshop addresses?


%% ─────────────────────────────────────────────────────────────
\section{Implications and Discussion Points}
%% ─────────────────────────────────────────────────────────────

What should the community do with your position? What open questions
or tensions does it surface? What would you like to discuss or debate
at the workshop?

Briefly sketch what a ``responsible practice'' or code of conduct
might look like in light of your argument.


%% ─────────────────────────────────────────────────────────────
%% AI Use Statement (required)
%% ─────────────────────────────────────────────────────────────
\begin{acks}
  \textbf{AI Use Statement:} [Required.] Describe how, if at all,
  you used Generative AI tools in writing this paper. Be specific:
  which tools, for which tasks (brainstorming, drafting, editing,
  figure generation, literature search, etc.), and what oversight
  you exercised. If you did not use AI tools, state that explicitly.

  Example: ``We used [Tool] for [task]. All claims, analysis, and
  framing are our own. We reviewed all AI-generated text and revised
  it substantially.''

  %% Additional acknowledgments (optional):
  % This work was supported by ...
\end{acks}

%% ─────────────────────────────────────────────────────────────
\bibliographystyle{ACM-Reference-Format}
\bibliography{references}
%% ─────────────────────────────────────────────────────────────

\end{document}
